\section{Experiments}
\label{sec:atm}

The experiments test whether post-recall failures are structured, predictable, and reducible.
We do not use the fact that aggregation is harder than lookup as a contribution.
Lookup queries serve as a sanity check: under oracle evidence, they should be easy.
The main evidence comes from three questions: whether evaluation complexity predicts failure, whether view insufficiency predicts error, and whether harnessed evaluation reduces semantic invalidity beyond prompting-only methods.

\subsection{ATM-Bench Diagnostic Annotation}
\label{subsec:benchmark-construction}

We use ATM-Bench as the base testbed because it evaluates long-term personalized memory QA over multimodal personal evidence, including images, videos, emails, metadata, and structured memory fields \cite{mei2026atm}.
To make post-recall failures measurable, we augment each analyzed instance with diagnostic annotations:
\begin{equation}
(\mathcal A, q, \mathcal A_q^\star, P_q, \mathcal Z_q(\mathcal A), \mathcal R_q, \mathcal L_q).
\label{eq:diagnostic-instance}
\end{equation}
Here $\mathcal A$ is the evidence structure, $q$ is the query, $\mathcal A_q^\star$ is the sufficient evidence substructure, $P_q$ is the gold evaluation program, $\mathcal Z_q(\mathcal A)$ is the semantic answer set, $\mathcal R_q$ is the set of required variables and relations, and $\mathcal L_q$ contains error labels for analysis.

The gold program $P_q$ specifies the operators needed to answer the query.
Operators include field selection, source selection, relation binding, joining, latest-record selection, override resolution, counting, summing, grouping, contradiction checking, and abstention under missingness.
The program is used for evaluation and diagnosis; it is not necessarily provided to the model.

The required-variable set $\mathcal R_q$ specifies the fields, timestamps, statuses, source identifiers, relation edges, operands, and state indicators required by $P_q$.
This enables view-level analysis.
A view is marked insufficient for $q$ if it fails to expose elements of $\mathcal R_q$ in a usable form.
We define the empirical view insufficiency rate of a view family $V$ as
\begin{equation}
\begin{aligned}
\mathrm{VIR}(V)
=
\Pr_{q,\mathcal A}
\big[
&V_q \text{ omits a required}\\
&\text{variable or relation from } \mathcal R_q
\big].
\end{aligned}
\label{eq:vir}
\end{equation}
This avoids treating flat text, tables, or state views as globally sufficient or insufficient.
Sufficiency is measured at the query-view level.

\subsection{Main Hypotheses}
\label{subsec:hypotheses}

\paragraph{H1: Evaluation complexity predicts post-recall failure.}
Under oracle evidence, semantic error should increase with binding count, operator count, override depth, conflict degree, and missingness level.
This tests the dynamic propagation prediction from Section~\ref{subsec:dynamic-binding}.

\paragraph{H2: View insufficiency predicts error.}
Query-view pairs with higher view insufficiency should have higher post-recall error.
Structured views and state views should reduce errors when they expose the variables and relations required by the gold program.

\paragraph{H3: Harnessed evaluation reduces semantic invalidity.}
Operator and verifier harnesses should reduce evidence-invalid answers more reliably than direct answering, chain-of-thought, or scratchpad prompting, especially for aggregation, temporal, conflict, and abstention queries.

The lookup-versus-aggregation comparison is reported only as a diagnostic.
It verifies that the oracle-evidence protocol is not broken, but it is not used as a main claim.

\subsection{Experimental Conditions}
\label{subsec:experimental-conditions}

We vary three factors: access, view, and intervention.

\paragraph{Access.}
Access conditions separate retrieval failure from post-recall failure.
In the retrieved condition, the model receives evidence from a standard retrieval pipeline.
In the oracle condition, it receives $\mathcal A_q^\star$.
In the oracle-plus-distractors condition, it receives $\mathcal A_q^\star$ together with irrelevant, stale, or confounding evidence.
The key post-recall analyses are performed under oracle and oracle-plus-distractors conditions.

\paragraph{View.}
View conditions test whether the recalled evidence is exposed in an evaluable form.
The flat view concatenates evidence as natural-language snippets.
The source-indexed view exposes item IDs, source IDs, timestamps, statuses, and field names.
The relation-explicit view exposes typed edges such as $\mathsf{newer\text{-}than}$, $\mathsf{overrides}$, $\mathsf{same\text{-}event}$, and $\mathsf{contradicts}$.
The state view exposes allowed decision-state variables, including candidate records, typed operands, source priorities, missingness flags, conflict flags, and relation labels.

Oracle-state views are controlled to avoid answer leakage.
They may expose operands and state variables required by the gold program, but they may not expose the final answer, the final natural-language conclusion, or the output of the target operator.
For lookup queries where the minimal state is identical to the answer, oracle-state views are excluded from the main comparison or reported only as an upper bound.

\paragraph{Intervention.}
Intervention conditions compare proposal shaping and evaluation-centric computation.
The proposal-shaping baselines are direct answering, chain-of-thought, and structured scratchpad.
The evaluation-centric interventions are operator harnesses and verifier harnesses.
Operator harnesses execute deterministic operations over extracted fields.
Verifier harnesses check whether candidate answers satisfy evidence-defined constraints.

\subsection{Complexity Controls}
\label{subsec:complexity-controls}

We vary evaluation complexity along five axes.
Binding count $k$ is the number of variables that must be simultaneously bound.
Operator count $o$ is the number of discrete operations in the gold program.
Override depth $c$ is the length of stale-to-current update chains.
Conflict degree $m$ is the number of mutually inconsistent or priority-ranked evidence items.
Missingness level measures how much of the answer-determining state is absent or underdetermined, requiring abstention.

Semantic complexity and context length are naturally coupled: more variables often require more evidence.
We therefore use two protocols.
In the natural-scaling protocol, evidence size is allowed to grow with complexity.
This reflects realistic memory use.
In the padded-control protocol, total context length and distractor count are approximately matched across complexity levels by adding irrelevant or template filler evidence.
A complexity effect is considered robust when it appears under both protocols.
If it appears only under natural scaling, we report it as a combined complexity-and-context effect.

\subsection{Models and Scope of Claims}
\label{subsec:models-scope}

We evaluate models in tiers.
Smaller open-weight models are used for large-scale sweeps, complexity curves, and error taxonomy.
Stronger open-weight or mid-scale models are used to test whether the trends survive increased capability.
Frontier closed models are used for smaller confirmation runs.

Claims are scoped by model tier.
A failure mode that appears only in smaller models is reported as scale-dependent.
A failure mode that persists across stronger models under oracle evidence is treated as stronger evidence for a general post-recall evaluation bottleneck.

\subsection{Metrics}
\label{subsec:metrics}

The primary metric is semantic accuracy:
\begin{equation}
\eta(y)\in\mathcal Z_q(\mathcal A).
\label{eq:semantic-accuracy}
\end{equation}
We also report semantic invalid rate, unsupported answer rate, binding accuracy, operator accuracy, abstention accuracy, and conflict-resolution accuracy.
For view analysis, we report $\mathrm{VIR}$ and the correlation between view insufficiency and semantic error.
For intervention analysis, we report verifier precision, verifier coverage, rejection rate, abstention rate, tool-call cost, token cost, and latency.

We report two gaps.
The oracle-recall gap measures the improvement from retrieved evidence to oracle evidence.
The post-recall gap measures the remaining difference between oracle evidence with a flat view and oracle evidence with state views or harnessed evaluation.

\subsection{Analysis Plan}
\label{subsec:analysis-plan}

The first analysis tests complexity scaling under oracle evidence.
We plot semantic accuracy and invalid rate against $k,o,c,m$, and missingness level.
This directly tests whether post-recall failure follows the complexity axes predicted by the dynamic propagation bound.

The second analysis tests view insufficiency.
We measure whether errors concentrate in query-view pairs where required fields, relations, or state indicators are absent.
We compare flat, source-indexed, relation-explicit, and state views under the same oracle evidence.
The key test is whether performance gains track reductions in $\mathrm{VIR}$, rather than merely changing the prompt format.

The third analysis compares prompting-only and harnessed evaluation.
We test whether chain-of-thought and scratchpads reduce errors by improving proposal quality, and whether operator or verifier harnesses further reduce semantic invalidity by executing or checking evidence-defined operations.

The final analysis assigns oracle-evidence failures to error categories: view error, binding error, temporal or state error, operator error, conflict error, abstention error, unsupported generation, and decoding error.
The purpose is not to show that hard queries are hard, but to show that post-recall errors are structured and can be predicted from the required evaluation program.

Together, these experiments determine the strength of the contribution.
If oracle evidence removes most errors, the bottleneck is mainly access.
If errors persist and scale with evaluation complexity, the results support a post-recall evaluation bottleneck.
If errors correlate with view insufficiency and decrease under state views, the results support decision-sufficient view as a mechanism.
If harnesses reduce semantic invalidity beyond prompting-only methods, the results support evaluation-centric memory reasoning rather than retrieval-only memory improvement.

\begin{table}[t]
\centering
\small
\begin{tabular}{lcccc}
\toprule
Condition & Access & Selection & View & Score \\
\midrule
Retrieval SGM & system & system & SGM & -- \\
NIAH SGM & oracle pool & system & SGM & -- \\
Oracle SGM & oracle & oracle & SGM & -- \\
Oracle Raw & oracle & oracle & raw & -- \\
Role View & oracle & oracle & role & -- \\
\bottomrule
\end{tabular}
\caption{Planned ATM-Bench diagnostic conditions. Scores will be filled after experiments.}
\label{tab:atm-diagnostic}
\end{table}

